\section{Game of Passwords}

In this section we outline a password gaming. The minimum number of questions to ask in order to guess their password with a high $\mathbb{P}$ in $x$ guesses.

Like the wordle game, what information about the target password would reveal useful info and therefore reduce the search space?

\begin{itemize}
    \item Length
    \item Structure
    \item Character placements
    \item Language
    \item Personal information (PII), e.g. date of birth.
\end{itemize}

Our assumptions are: 1) the passwords are user generated, 2) we are not allowed to ask what the password is directly.

\section{Markov Model}
An alphabetical password generated by a human, even if it is not a dictionary word, is unlikely to be uniformly distributed in the space of alphabet sequences. In fact, if asked to pick a sequence of characters at random, it is likely that an English-speaking user will generate a sequence in which each character is roughly equidistributed with the frequency of its occurrence in English text. Analysis of our password database reveals a significant number of alphabetical passwords which are neither dictionary words nor random sequences (we used the openwall.com dictionary which contains about 4 million words \cite{openwall}).

Markov models are commonly used in natural language processing and are at the heart of speech recognition systems \cite{jurafsky}. A Markov model defines a probability distribution over sequences of symbols. In other words, it allows sampling character sequences that have certain properties. In fact, Markov models have been used before in the context of passwords. Subsequent to this work, we learned of the \emph{Extensible Multilingual Password Generator} software, which was written by Jon Callas in 1991 and used precisely this technique to generate passwords for users. (We are also told that a similar method for generating "random, yet pronounceable" passwords was used at the Los Alamos National Laboratory in the late 1980s.) It should come as no surprise, then, that Markov modeling is very effective at guessing passwords generated by users.

In a zero-order Markov model, each character is generated according to the underlying probability distribution and independently of the previously generated characters. Mathematically, the probability of a sequence $S = c_1 c_2 \ldots c_n$ of $n$ characters is given by:

\[
P(S) = \prod_{i=1}^{n} P(c_i),
\]

where $P(c_i)$ is the probability of character $c_i$ occurring.

In a first-order Markov model, each character depends on the immediately preceding character. Each bigram (ordered pair of characters) has an associated probability, and the sequence is generated by considering these conditional probabilities. Mathematically, the probability of the sequence $S$ is:

\[
P(S) = P(c_1) \prod_{i=2}^{n} P(c_i \mid c_{i-1}),
\]

where $P(c_i \mid c_{i-1})$ is the conditional probability of character $c_i$ given that the previous character is $c_{i-1}$.

In general, for a $k$-th order Markov model, the probability of the sequence $S$ is:

\[
P(S) = \prod_{i=1}^{k} P(c_i) \prod_{i=k+1}^{n} P\left(c_i \,\middle|\, c_{i-k}, c_{i-k+1}, \ldots, c_{i-1}\right),
\]

where $P\left(c_i \,\middle|\, c_{i-k}, c_{i-k+1}, \ldots, c_{i-1}\right)$ is the probability of character $c_i$ given the preceding $k$ characters.

By utilizing higher-order Markov models, we can capture longer dependencies and more accurately model the statistical properties of human-generated passwords. This approach allows us to generate password candidates that are more representative of those chosen by users, improving the effectiveness of password guessing and security assessments.

To construct these Markov models, we estimate the probabilities from a training dataset, such as a corpus of English text or a collection of known passwords. The character probabilities $P(c_i)$ and the conditional probabilities $P(c_i \mid c_{i-1})$ (or higher-order conditional probabilities) are calculated by counting the occurrences in the dataset:

\[
P(c_i) = \frac{\text{Count}(c_i)}{\sum_{c \in \mathcal{C}} \text{Count}(c)},
\]

\[
P(c_i \mid c_{i-1}) = \frac{\text{Count}(c_{i-1} c_i)}{\sum_{c' \in \mathcal{C}} \text{Count}(c_{i-1} c')},
\]

where $\mathcal{C}$ is the set of all possible characters, $\text{Count}(c_i)$ is the number of times character $c_i$ appears, and $\text{Count}(c_{i-1} c_i)$ is the number of times the bigram $c_{i-1} c_i$ appears in the dataset.

These models not only aid in understanding the structure of user-generated passwords but also play a crucial role in designing more secure authentication systems by highlighting weaknesses in common password choices.